\begin{proposition}\hypertarget{tag-01FW}{}\label{cdp:st:D.1.11}
\StatementRecord{cdp:st:D.1.11}{proved}{01FW}{\SeriesRef{cdp:st:D.1.9}, \SeriesRef{cdp:st:D.1.7}, \SeriesRef{cdp:st:D.1.8}}

Let $f:Y\to\PP^1$ be a proper connected-fibre holomorphic map from
a connected compact complex manifold with $b_1(Y)=b_2(Y)=0$.
In every component of $\Pic(Y)$, the line
bundles which are non-torsion on every full fibre form a dense
$G_\delta$ set of full Lebesgue measure in its complex vector-space
coordinates. If evaluation of global sections of $R^1f_*\OO_Y$
has rank one, then $\Pic^0(Y)=\C$ and the exceptional set is
countable. For two-torus fibres this alternative holds whenever
the splitting has a negative summand.
In particular the general-line-bundle conclusion of CDP20
Proposition 5.6 holds with this precise quantifier, without the
nonnegative-summand assertion of Proposition 5.5.
\end{proposition}


\begin{proof}

\begin{enumerate}[label={$\langle1\rangle$\arabic*.},leftmargin=*]
\item Let $E=R^1f_*\OO_Y$, a vector bundle by \cite[Theorem~2.2, pp.~4--5]{Viaclovsky2026Combs}, and use
$\gamma,\beta$ from \ref{cdp:st:D.1.9}. The section of $E$ corresponding to
$\beta$ has no zero. In fact its value maps naturally to the
nonzero class $\iota_c(\gamma_c)$ in the cohomology of every full
fibre; Lemmas \ref{cdp:st:D.1.7} and \ref{cdp:st:D.1.8} prove nonvanishing. Write
$E=\bigoplus_j\OO(a_j)$ and let $k$ count the nonnegative degrees.
Evaluation has constant rank $k$ onto their direct-summand subbundle
$E_+$. Nonvanishing of $\beta$ gives $k\geq1$. When $k=1$, its
sole nonnegative degree is zero, since a positive-degree line bundle
has no nowhere-zero section. In that case $H^0(E)=\C\beta$.

\item Suppose first $k\geq2$, and put $V=H^1(Y,\OO_Y)=H^0(E)$ and
$N=\dim_{\C}V$. On a simply connected disc in the smooth-fibre
locus where degree-one base change holds, trivialize the integral
local system. For a flat integral
class $\lambda$, the equation
\[
z|_{F_c}=\iota_c(\lambda),\qquad (z,c)\in V\times\Delta
\]
defines an analytic set. The parameter values where
$\iota_c(\lambda)\in(E_+)_c$ form an analytic subset of the disc;
over it the fibres are affine spaces of dimension $N-k$.
The incidence set thus has complex dimension at most
$N-k+1\leq N-1$. Its projection to $V$ has Lebesgue
measure zero. It is also meagre: stratify and exhaust that set
by countably many compact sets; the images are closed sets with
empty interior, by their real dimension at most $2N-2$.
Take countably many discs, integral classes and positive integer
denominators to include every torsion condition.

\item There are finitely many excluded fibres, including the exceptional
values of degree-one base change. At each, the natural map
$E\otimes\C(c)\to H^1(Y_c,\OO_{Y_c})$ is injective. This follows
directly from the long exact sequence for multiplication by a local
base parameter: its cokernel is the parameter-torsion in
$R^2f_*\OO_Y$. The restriction $V\to H^1(Y_c,\OO_{Y_c})$
therefore has rank $k\geq2$. The torsion conditions are contained in
countably many affine subspaces of complex codimension at least
two. Combining these with the preceding projections gives a
meagre null set. The complement is dense and $G_\delta$, as claimed.

\item It remains to prove the case $k=1$.
Now $V=\C\beta$. On a simply connected disc in the same open locus,
the projection of a flat integral class $\lambda$ to
$E/(\OO\beta)$ is holomorphic. Unless it vanishes identically,
its zero set is discrete. Such zeroes give only countably many
parameters $z$ for which $\iota_c(\lambda)=z\beta_c$.
We prove that an identically vanishing projection gives a constant
parameter, including after analytic continuation.

\item Pass to the universal cover of the smooth locus, on which
$H^1(F;\Z)$ is marked. The identity theorem continues the relation
$\iota_c(\lambda)\in\C\iota_c(\gamma)$ everywhere on this cover.
If $\lambda$ and $\gamma$ are real-dependent, their ratio is a
constant rational number. Otherwise let
$W_{\Q}=\Q\gamma+\Q\lambda$. The coefficient map
$H^1(F;\R)\to H^1(F,\OO_F)$ is real-linearly injective by \ref{cdp:st:D.1.7}.
The images of the two real-independent classes span the entire
real two-plane underlying $\C\beta_c$, so $W_{\R}$ is precisely
its inverse image. Analytic continuation around a base loop
sends $\lambda$ to a class whose coefficient image again belongs
to this line. Hence it preserves $W_{\Q}$ and fixes $\gamma$.

\item Choose an integral basis $(\gamma_0,\lambda_0)$ of
$W_{\Q}\cap H^1(F;\Z)$ with $\gamma$ a positive integral multiple
of $\gamma_0$. The induced monodromy on this lattice fixes
$\gamma_0$ and preserves the orientation of $W_{\R}$: the complex
structure induced on that plane by the coefficient map varies
continuously along the loop, so its complex orientation cannot change.
It therefore sends $\lambda_0$ to $\lambda_0+k\gamma_0$ for an
integer $k$. Put
\[
\tau(c)=\frac{\iota_c(\lambda_0)}{\iota_c(\gamma_0)}.
\]
This holomorphic function on the universal cover never takes a
real value, by real injectivity of the coefficient map, and its
imaginary part has constant sign. Monodromy changes $\tau$ by
integers. Accordingly either $\exp(2\pi i\tau)$ or
$\exp(-2\pi i\tau)$ descends to a holomorphic function of absolute
value less than one on $\PP^1$ minus finitely many points. The
removable singularity theorem extends it to $\PP^1$, where it is
constant. Its logarithmic derivative gives $d\tau=0$. Thus the
original ratio is constant as asserted.

\item There are countably many integral classes on a countable disc
cover. The isolated-zero alternatives and the constant alternatives
just proved give a countable exceptional set on smooth fibres.
Each exceptional full fibre contributes another countable set,
because $z\mapsto z\beta_c$ is injective and its intersection with
the rational integral coefficient image is countable. This proves
the assertion in $\Pic^0(Y)$.

\item Finally for any component $A\Pic^0(Y)$, the map taking the $N$th
power, for an integer $N$ killing its component class, is the affine
isomorphism $z\mapsto b+Nz$ of the complex parameter space used
in \ref{cdp:st:D.1.9}. It preserves the torsion condition in every fibre, and
preserves countability, meagreness and measure-zero sets. This
proves all the componentwise assertions.

\end{enumerate}
\end{proof}